\section{Related Work}
\label{sec:related}

Research relevant to adversarial cybersickness detection spans four streams: AI-based cybersickness detection, adversarial VR threat modelling, explainable AI (XAI) for safety-critical detection, and value-based reinforcement learning for detection. We summarise each and the gap that motivates E-DDQN. Accuracy figures quoted below are those reported by the cited authors on their own datasets and evaluation protocols; they are not comparable with each other or with ours, and most do not state whether training and test data share participants.

\subsection{AI-Based Cybersickness Detection}
Questionnaires such as the Simulator Sickness Questionnaire~\cite{kennedy1993ssq} measure discomfort after exposure and cannot capture onset during immersion~\cite{caserman2021cybersickness}, which motivates sensor-driven detection. Dennison et al.~\cite{dennison2016physio} related physiological signals (e.g., breathing and stomach activity) to cybersickness reports. Islam et al. detected and forecast severity from heart rate, breathing rate, heart-rate variability and skin conductance with deep neural classifiers~\cite{islam2020automatic}, and later fused eye- and head-tracking data from integrated HMD sensors (87.77\,\% accuracy)~\cite{islam2021fusion}. Liu et al. assessed cybersickness from EEG with a CNN-ECA-LSTM network~\cite{liu2024cnneca}; Yalcin et al. combined bidirectional LSTMs with conditional-GAN augmentation of consumer-grade physiological data (91.1\,\% F1)~\cite{yalcin2024automatic}; Jeong and Han used transformers on multimodal time series~\cite{jeong2024precyse}. Using gameplay data, Porcino et al. collected UFFCSData and used symbolic machine learning to identify likely causes such as rotation and acceleration~\cite{porcino2022identifying}. A systematic review notes that most detectors classify the current state rather than anticipating onset~\cite{yang2022mlreview}. None of these works evaluates detection under deliberate manipulation, and few report participant-independent evaluation.

\subsection{Adversarial VR Threats}
VR mediates perception and balance and therefore exposes an attack surface distinct from conventional software~\cite{odeleye2023virtually,giaretta2025security}. Tseng et al. showed how perceptual manipulations (e.g., redirected walking and haptic retargeting) can be repurposed to cause physical harm~\cite{tseng2022dark}, and Cheng et al. studied users' reactions to perceptual-manipulation attacks in mixed reality~\cite{cheng2023perceptual}. Closer to our threat model, Odeleye et al. described frame-rate manipulation attacks that can induce VR sickness and detected them early with an Isolation Forest~\cite{odeleye2021framerate}, and Kundu et al. introduced \emph{cybersickness attacks}---adversarial inputs that fool deep cybersickness detectors so that mitigation is not triggered---together with an XAI-guided defence~\cite{kundu2025securing}. Network-level anomaly detection for VR learning environments addresses denial-of-service and unauthorised access rather than perception-layer manipulation~\cite{valluripally2023detection}.

\subsection{Explainable AI for Safety-Critical Detection}
SHAP~\cite{lundberg2017shap} and LIME~\cite{ribeiro2016lime} provide per-decision attributions. In cybersickness detection, TruVR used explainable boosting machines (99.75\,\% and 94.10\,\% accuracy on physiological and gameplay data)~\cite{kundu2022truvr}, LiteVR used SHAP-guided feature reduction for lightweight models~\cite{kundu2023litevr}, and VR-LENS used XAI-guided feature reduction to shorten training and inference on a mobile processor~\cite{kundu2023vrlens}. Kundu et al. further used SHAP ``signatures'' to separate normal from adversarial detector inputs~\cite{kundu2025securing}. Post hoc attribution computed off the inference path adds no latency to the detection decision itself.

\subsection{Value-Based Reinforcement Learning for Detection}
Deep Q-networks~\cite{mnih2015human} suffer from overestimation, which Double DQN reduces by decoupling action selection and evaluation~\cite{vanhasselt2016double}; dueling heads separate state value from action advantage~\cite{wang2016dueling}. Classification has been cast as a sequential decision process solved with DQN, with larger rewards for minority-class samples, by Lin et al.~\cite{lin2020imbalanced}; asymmetric rewards for detection are therefore not new in themselves. Their relation to classical cost-sensitive learning~\cite{elkan2001cost} is, however, rarely made explicit, and safe-RL research shows that the choice of penalties shapes how risk-averse a learned policy is~\cite{garcia2015safe}. To our knowledge, asymmetric-reward value learning has not been evaluated for cybersickness detection under a participant-independent protocol.

\subsection{Identified Research Gap}
Table~\ref{tab:gaps} summarises the gaps. Three are addressed directly in this work: a safety-prioritised (5:1) decision objective whose relation to cost-sensitive thresholding is made explicit; a leakage-free, participant-conservative evaluation that compares the RL detector with matched supervised alternatives; and a reproducible stress test of the detector under bounded manipulation of the telemetry channels that an adversary could influence. Physically validated rendering attacks remain outside what any public dataset currently supports.

\begin{table*}[t]
\centering
\caption{Identified research gaps and how this revision addresses them.}
\label{tab:gaps}
\footnotesize
\begin{tabular}{p{0.2\textwidth}p{0.27\textwidth}p{0.45\textwidth}}
\toprule
Gap & Evidence & This work \\
\midrule
Detectors are rarely evaluated under deliberate manipulation & Only \cite{kundu2025securing,odeleye2021framerate} consider attacks & Bounded, reproducible telemetry-perturbation stress tests (frame pacing, camera motion, locomotion) with evasion and false-alarm worst cases; physical attacks remain future work \\
Symmetric objectives ignore unequal error costs & \cite{mnih2015human,vanhasselt2016double} & 5:1 asymmetric reward; proof that it equals the 5:1 cost-sensitive Bayes threshold (Eq.~\ref{eq:margin}); evaluated against cost-sensitive supervised models at matched operating points \\
Participant leakage and optimistic splits & \cite{kapoor2023leakage,kaufman2012leakage,roberts2017cv} & Session-linkage groups, nested grouped evaluation, train-only preprocessing, frozen protocol \\
Explanations that do not target the decision & SHAP on class scores & SHAP on the decision margin $\Delta Q$ \\
\bottomrule
\end{tabular}
\end{table*}
